\section{Full-germ correspondence calculus}\label{sec:calculus}

The calculations in this section are algebraic.  We give a group-factor
version of the projective induction and cocycle-corrected double-coset
calculus in Vaes's work on generalized Bernoulli crossed products
\cite[Theorem~2.2 and Section~6.2, especially Definition~6.8 and
Theorem~6.9]{Vaes2008}.  We prove the formulas here, including the
full-stabilizer argument for every bounded intertwiner and every bimodule
summand.  Until we apply
Proposition~\ref{prop:recovery-exhaustion}, the groups need only be countable
and ICC; the geometric hypotheses of Definition~\ref{def:class} are not
used.  We first construct the arithmetic models with their exact scalar
transports, and then determine their intertwiners, summands, conjugates,
and relative tensor products.  This order also verifies that every model
appearing in the exhaustion is a normal bifinite correspondence.

Write
\[
 M=L_\mu(\Gamma),\qquad N=L_\omega(\Lambda),\qquad
 D=\Gamma\times\Lambda,
 \qquad
 \alpha((g,h),(g',h'))=\mu(g,g')\overline{\omega(h,h')}.
\]
On an $M$--$N$ correspondence the joint group action is
\begin{equation}\label{eq:calculus-joint-action}
 U_{(g,h)}\xi=u_g\xi v_h^*.
\end{equation}
It has multiplier $\alpha$.  Indeed, its rightmost product is
$v_{h'}^*v_h^*=(v_hv_{h'})^*$, so
$U_{(g,h)}U_{(g',h')}=
\mu(g,g')\overline{\omega(h,h')}U_{(gg',hh')}$.
In particular the two coordinate actions commute exactly.

\subsection{Germs and their full stabilizers}

A germ from $\Gamma$ to $\Lambda$ is an isomorphism
$\delta:A\to B$ between finite-index subgroups, modulo agreement after
restriction to a common finite-index subgroup.  Composition means
composition on a finite-index domain where both representatives are
defined.  For example, if $\varepsilon:C\to F$ is a representative from
$\Lambda$ to a third group, then $\varepsilon\delta$ is defined on
$\delta^{-1}(B\cap C)$ and has image $\varepsilon(B\cap C)$.
Both have finite index in the relevant endpoint group.  Restricting
representatives further does not change the germ, and a common further
restriction verifies associativity.  Inversion is represented by the
inverse subgroup isomorphism.

For a germ $d$ from $\Gamma$ to $\Lambda$, let
\[
 (g,h)d=\Ad(h)\,d\,\Ad(g^{-1}),\qquad
 C_d=\{(g,h)\in D:(g,h)d=d\}.
\]
If $\delta:A\to B$ represents $d$, denote its graph by
$C_\delta=\{(a,\delta(a)):a\in A\}$.
We distinguish this graph from the full stabilizer $C_d$.

\begin{lemma}[Full stabilizers and the finite-orbit test]
\label{lem:calculus-germ-test}
The maps $g\mapsto[\Ad(g)]$ and $h\mapsto[\Ad(h)]$ are injective.
The group $C_d$ is the graph of an isomorphism
$\delta_d:A_d\to B_d$ between finite-index subgroups.  It contains
$C_\delta$ at finite index, and $\delta_d$ extends $\delta$.
If a finite-index subgroup of $C_d$ fixes another germ $x$ from
$\Gamma$ to $\Lambda$, then $x=d$.  Equivalently, the only finite
$C_d$-orbit in the space of these germs is $\{d\}$.
The same conclusion holds with $C_d$ replaced by $C_\delta$.
\end{lemma}

\begin{proof}
The centralizer of a finite-index subgroup of an ICC group is trivial:
an element of that centralizer has a centralizer of finite index and
therefore a finite conjugacy class.  This proves injectivity of the
inner-germ maps.

The stabilizer condition can be written
\begin{equation}\label{eq:calculus-stabilizer-relation}
 \Ad(h)=d\Ad(g)d^{-1}.
\end{equation}
If either coordinate of an element of $C_d$ is the identity, injectivity
of the inner-germ maps forces the other coordinate to be the identity.
Thus both coordinate projections on $C_d$ are injective.  Their images
$A_d$ and $B_d$ form the domain and range of a group isomorphism
$\delta_d$.  The inclusion $C_\delta\subset C_d$ follows by evaluating
Equation~\eqref{eq:calculus-stabilizer-relation} on a common finite-index
domain.  Consequently $A\subset A_d$, $B\subset B_d$, and
$\delta_d|_A=\delta$.  Both projected groups have finite index, and
$[C_d:C_\delta]=[A_d:A]<\infty$.

Suppose $C_0\leq C_d$ has finite index and fixes $x$.  Its first
projection $A_0$ has finite index in $\Gamma$.  For every
$(a,b)\in C_0$ the two stabilizer equations give
\begin{equation}\label{eq:calculus-finite-orbit-test}
 x\Ad(a)x^{-1}=\Ad(b)=d\Ad(a)d^{-1}.
\end{equation}
Choose representatives $\chi$ and $\delta$ for $x$ and $d$, and restrict
$a$ to $A_0\cap\operatorname{dom}(\chi)\cap
\operatorname{dom}(\delta)$.  On this finite-index subgroup the outer
expressions in Equation~\eqref{eq:calculus-finite-orbit-test} are the
inner germs of $\chi(a)$ and $\delta(a)$.  Injectivity in $\Lambda$
gives $\chi(a)=\delta(a)$, whence $x=d$.  A finite orbit has a
finite-index point stabilizer, so this proves the equivalent orbit
statement.  Finally, a finite-index subgroup of $C_\delta$ also has
finite index in $C_d$.
\end{proof}

The finite-orbit conclusion will make germ labels intrinsic to the
induced models.  Cosets of $C_d$ correspond to actual germs, whereas
several cosets of a presenting graph may represent the same germ.
After grouping those cosets, a bounded intertwiner will have matrix
blocks between the resulting finite-dimensional fibres.  Equivariance makes a block's
norm constant along a stabilizer orbit.  The lemma will therefore force
such a map to preserve each germ label, once square summability of its
columns is used.

\subsection{Projective induction and the arithmetic model}

Projective induction, induction in stages, and double-coset tensor
decompositions belong to Mackey's representation-theoretic framework
\cite[Theorems~4.1, 4.2, and~4.6]{Mackey1958}. The following coordinates fix the
multiplier throughout and prove every scalar identity used here.

Let $C\leq D$ be a subgroup and let $\gamma$ be a normalized unitary
projective representation on a finite-dimensional space $V$, with
multiplier $\alpha|_{C\times C}$.  Write $X=D/C$ for left cosets.
Choose representatives $r_x\in D$, with $r_C=e$, and give each $x\in X$
a copy $V_x$ of $V$.  Set
$\mathcal H(C,\gamma)=\bigoplus_{x\in X}V_x$, with the counting
Hilbert norm.  For $s\in D$ put
\begin{equation}\label{eq:input-A1}
 \begin{split}
 c(s,x)&=r_{sx}^{-1}sr_x,\\
 U^X_s(x)&=
 \frac{\alpha(s,r_x)}{\alpha(r_{sx},c(s,x))}\,
 \gamma(c(s,x)):V_x\longrightarrow V_{sx}.
 \end{split}
\end{equation}

\begin{lemma}[Exact projective induction]
\label{lem:calculus-induction}
Equation~\eqref{eq:input-A1} defines a unitary projective representation
of $D$ with multiplier $\alpha$.  Changing the coset representatives
gives a unitarily equivalent bundle by unitary identifications of its
individual fibres.  If $C$ is the graph of a finite-index subgroup
isomorphism $A\to B$, the bundle has commuting normal $M$ and right
$N$ actions, and is bifinite.
\end{lemma}

\begin{proof}
The correction $c(s,x)$ belongs to $C$, and
\[
 c(s,tx)c(t,x)=c(st,x).
\]
Using the cocycle identity for $\alpha$ and the projective multiplication
of $\gamma$ in Equation~\eqref{eq:input-A1} gives
\begin{equation}\label{eq:calculus-transport-identity}
 U^X_s(tx)U^X_t(x)=\alpha(s,t)U^X_{st}(x).
\end{equation}
One can also check all the scalar factors at once in the central
extension with multiplication
$(z,s)(w,t)=(zw\alpha(s,t),st)$: the representation of its preimage
of $C$ is $(z,c)\mapsto z\gamma(c)$, and the equality
$sr_x=r_{sx}c(s,x)$ gives exactly the ratio in
Equation~\eqref{eq:input-A1}.  Thus this argument retains, rather than
discards, the fixed multiplier.

For explicit independence of representatives, suppose
$r'_x=r_xk_x$, where $k_x\in C$.  The identification from the new
coordinate copy of $V_x$ to the old one is
\[
 F_x=\alpha(r_x,k_x)^{-1}\gamma(k_x).
\]
Indeed identifying a fibre by $U_{r'_x}$ instead of $U_{r_x}$ has
precisely this effect.  Equation~\eqref{eq:calculus-transport-identity}
then gives $F_{sx}U'{}^X_s(x)=U^X_s(x)F_x$.

Now assume that $C$ is a graph.  Both coordinate groups act freely on
$X$.  For example, if $(g,e)r_xC=r_xC$, then
$r_x^{-1}(g,e)r_x\in C$ has second coordinate $e$, so $g=e$.
On an orbit of the first coordinate group choose a fibre $V_{x_0}$.
The unitary identification
\[
 \ell^2(\Gamma)\otimes V_{x_0}\longrightarrow
 \bigoplus_{g\in\Gamma}V_{(g,e)x_0},
 \qquad \delta_g\otimes\xi\longmapsto U^X_{(g,e)}\xi
\]
intertwines its action with $\lambda_\mu\otimes1$, since
$U_{(a,e)}U_{(g,e)}=\mu(a,g)U_{(ag,e)}$.  Hence it extends normally
to $L_\mu(\Gamma)$.  The same argument for the second coordinate
gives copies of $\lambda_{\bar\omega}$, which are the standard
normal representations of the generators acting by right adjoints
$\xi\mapsto\xi v_h^*$.  The two actions commute, because the product
cocycle has trivial mixed-coordinate factors.

The first coordinate orbits are indexed by $\Lambda/B$, and the
second coordinate orbits by $\Gamma/A$.  They are finite in number.
Each is a regular module with multiplicity $\dim_{\mathbb C}V$,
proving bifiniteness as well as the dimension formulas below.
\end{proof}

We next identify this construction with the Connes-product model in
Theorem~\ref{thm:correspondences}.  This also supplies a direct normality
check of the subgroup action used to define that product.

\begin{lemma}[The graph-bundle unitary]
\label{lem:calculus-graph-model}
Let $\delta:A\to B$ be a finite-index subgroup isomorphism and let
$\sigma:A\to\mathcal U(V)$ have multiplier
$\beta=\mu|_{A\times A}/\delta^*\omega$.  Define
$\gamma(a,\delta(a))=\sigma(a)$.  Then
\[
 E(A,B,\delta,\sigma)=
 L^2(M)\mathbin{\bar\otimes}_{L_{\mu|A}(A)}
       (V\otimes L^2(N))
 \cong\mathcal H(C_\delta,\gamma).
\]
The unitary sends the vector $\xi$ in the fibre labelled by $x$ to
$U_{r_x}j\xi$, where
$j\xi=\widehat1\otimes(\xi\otimes\widehat1)$ in the displayed model.
\end{lemma}

\begin{proof}
On $V\otimes L^2(N)$ the subgroup action is
\[
 u_a(\xi\otimes\eta)=\sigma(a)\xi\otimes v_{\delta(a)}\eta.
\]
Its multiplier is $\beta(a,b)\omega(\delta(a),\delta(b))=\mu(a,b)$.
For normality, choose representatives $t$ for $B\backslash\Lambda$.
On the summand spanned by $\xi\otimes v_{\delta(a)}v_t$, the map
\begin{equation}\label{eq:calculus-absorption}
 \xi\otimes v_{\delta(a)}v_t
 \longmapsto\delta_a\otimes\sigma(a)^*\xi
 \quad\text{in }\ell^2(A)\otimes V
\end{equation}
is unitary and conjugates the subgroup action to
$\lambda_{\mu|A}\otimes1$.  To check the phase, applying $u_b$
before this map produces
\[
 \omega(\delta(b),\delta(a))
 \sigma(ba)^*\sigma(b)
 =\omega(\delta(b),\delta(a))\beta(b,a)\sigma(a)^*
 =\mu(b,a)\sigma(a)^*.
\]
Thus the left subgroup action is normal.  Its right $N$ action is
the standard action on the second factor and commutes with it.

Choose representatives $r$ for $\Gamma/A$.  As a right
$L_{\mu|A}(A)$ module, $L^2(M)$ has orthonormal free basis $u_r$;
the assertion follows directly from
$E_{L_{\mu|A}(A)}(u_r^*u_{r'})=\delta_{r,r'}1$.
The relative product therefore identifies with the orthogonal sum
\begin{equation}\label{eq:calculus-coset-model}
 E(A,B,\delta,\sigma)
 =\bigoplus_{r\in\Gamma/A}u_r\otimes(V\otimes L^2(N)).
\end{equation}
At its identity fibre the graph action is exactly
\[
 U_{(a,\delta(a))}j\xi
 =u_a\otimes(\xi\otimes v_{\delta(a)}^*)
 =j\sigma(a)\xi.
\]
For $g=ra$, balancing in Equation~\eqref{eq:calculus-coset-model}
gives the more explicit formula
\[
 U_{(g,h)}j\xi
 =\mu(r,a)^{-1}u_r\otimes
       \bigl(\sigma(a)\xi\otimes v_{\delta(a)}v_h^*\bigr).
\]
The last Fourier coordinate is $\delta(a)h^{-1}$, up to its displayed
unit scalar.  Thus two translated identity fibres are orthogonal
unless their labels in $D/C_\delta$ agree: equality of their first
cosets and last Fourier coordinates is exactly that coset equality.
Conversely these fibres exhaust every summand in
Equation~\eqref{eq:calculus-coset-model}.  They are an orthogonal,
transitive system with graph action $\gamma$ at the basepoint.
Finally
\[
 U_sU_{r_x}j
 =\frac{\alpha(s,r_x)}{\alpha(r_{sx},c(s,x))}
       U_{r_{sx}}j\gamma(c(s,x)),
\]
which proves that the stated unitary has exactly the transport
of Equation~\eqref{eq:input-A1}.
\end{proof}

For a full stabilizer we write
\[
 H(d,\gamma)=\mathcal H(C_d,\gamma).
\]
Its labels identify with actual germs in the orbit $Dd$.
There is also a canonical regrouping of an original graph model over
this full-germ orbit.  The map $D/C_\delta\to D/C_d$ has finite
fibres.  Its fibre above $C_d$ carries
\begin{equation}\label{eq:calculus-regrouping}
 \operatorname{Ind}^{C_d,\alpha}_{C_\delta}\gamma,
\end{equation}
where projective induction uses
Equation~\eqref{eq:input-A1} with $D$ replaced by $C_d$ and with the
restricted cocycle.  Indeed that fibre is the sum indexed by
$C_d/C_\delta$, and its transports are precisely that formula.
Regrouping the orthogonal sum consequently gives
\[
 \mathcal H(C_\delta,\gamma)
 \cong H\bigl(d,\operatorname{Ind}^{C_d,\alpha}_{C_\delta}\gamma\bigr).
\]
This description includes the phases of finite induction; it does not
replace the projective representation by an ordinary one.

Changing an orbit representative also retains a scalar transport.
If $d'=kd$, then $C_{d'}=kC_dk^{-1}$.  Identify the new base fibre
with the old one using $U_k$.  For $c'=kck^{-1}\in C_{d'}$ the
representation in these coordinates is
\begin{equation}\label{eq:calculus-representative-change}
 \gamma'(c')=
 \frac{\alpha(c',k)}{\alpha(k,c)}\,\gamma(c).
\end{equation}
Indeed $U_{c'}U_k=
\alpha(c',k)\alpha(k,c)^{-1}U_kU_c$.
This gives the exact representative-change equivalence, including its
restricted multiplier.  Two choices of $k$ leading to the same germ
give equivalent fibres by the stabilizer action itself.

\begin{proposition}[Dimensions of the actual orbit modules]
\label{prop:calculus-dimensions}
For a finite-dimensional $\alpha|_{C_d}$-projective representation
$\gamma$,
\begin{equation}\label{eq:calculus-full-dimensions}
 \dim_M H(d,\gamma)
   =\dim_{\mathbb C}(\gamma)[\Lambda:B_d],
 \qquad
 \dim H(d,\gamma)_N
   =\dim_{\mathbb C}(\gamma)[\Gamma:A_d].
\end{equation}
For the original model these formulas read
\[
 \dim_M E(A,B,\delta,\sigma)
   =\dim_{\mathbb C}(\sigma)[\Lambda:B],
 \qquad
 \dim E(A,B,\delta,\sigma)_N
   =\dim_{\mathbb C}(\sigma)[\Gamma:A].
\]
Zero fibres may be omitted.  Dimensions add under finite orthogonal
sums.
\end{proposition}

\begin{proof}
The free coordinate-orbit identifications in
Lemma~\ref{lem:calculus-induction} exhibit exactly the stated numbers
of copies of the standard left and right modules.  Each standard
module has dimension one for the normalized factor trace.
Lemma~\ref{lem:calculus-graph-model} gives the original-model formulas.
\end{proof}

\subsection{All bounded intertwiners and all summands}

\begin{theorem}[Full-germ intertwiner formula]
\label{thm:calculus-hom}
Choose finitely many distinct $D$-orbits of germs and a representative
$d$ in each orbit.  Let $\gamma_d,\tau_d$ be finite-dimensional
projective representations of $C_d$, with multiplier $\alpha|_{C_d}$;
either may be zero.  Then
\begin{equation}\label{eq:calculus-hom}
 \Hom_{M\text{--}N}
 \left(\bigoplus_dH(d,\gamma_d),\,
       \bigoplus_dH(d,\tau_d)\right)
 =\bigoplus_d\Hom_{C_d,\alpha}(\gamma_d,\tau_d).
\end{equation}
The equality is implemented by restriction to the representative
fibres and extension by transport.  Norm is the maximum of the fibre
norms, and adjoints and composition are fibrewise.  In particular all
closed bimodule summands are obtained by taking invariant subspaces
of these finite-dimensional fibre representations.
\end{theorem}

\begin{proof}
Realize both modules over their union of full-germ orbits, assigning a
zero fibre where necessary.  An intertwiner $T$ has finite-dimensional
matrix blocks $T_{yx}:V_x\to W_y$.  If $(e_i)$ is an orthonormal basis
of $V_x$, then
\begin{equation}\label{eq:calculus-hs-column}
 \sum_y\|T_{yx}\|_{\mathrm{HS}}^2
 =\sum_i\|Te_i\|^2
 \leq\dim(V_x)\|T\|^2<\infty.
\end{equation}
Equivariance identifies $T_{sy,sx}$ with
$U^\tau_s(y)T_{yx}U^\gamma_s(x)^*$.  Its Hilbert--Schmidt norm is
therefore unchanged.  For $s\in C_x$, Equation~\eqref{eq:calculus-hs-column}
implies that a nonzero block $T_{yx}$ can occur only when the
$C_x$-orbit of $y$ is finite.  By
Lemma~\ref{lem:calculus-germ-test}, this forces $y=x$.
Thus every intertwiner is diagonal in the germ labels, including
between initially different orbit bundles.

At a representative label $d$, equivariance is exactly the requirement
that $T_d$ intertwine the two $\alpha|_{C_d}$-projective
representations.  Conversely, such an operator extends by
\[
 T_{sd}=U^\tau_s(d)T_dU^\gamma_s(d)^*.
\]
This is independent of the choice of $s$: changing $s$ by an element
of $C_d$ uses the intertwining identity, and the identical scalar
multiplier on both sides cancels.  The extension is bounded with norm
$\|T_d\|$ on that orbit and commutes with every joint group generator.
Normality of the endpoint representations then makes it a bimodule
intertwiner.  This proves Equation~\eqref{eq:calculus-hom} and the
norm assertion.  The formulas for adjoint and composition follow
on each fibre.

A closed bimodule summand is the range of a projection in the bounded
bimodule commutant.  The preceding description identifies that
projection with a family of $C_d$-equivariant fibre projections.
Their ranges are precisely the asserted projective-representation
summands, and their transported ranges give the entire module summand.
\end{proof}

Consider any closed bimodule summand $K$ of a finite sum of the original
graph models. Regroup each model by Equation~\eqref{eq:calculus-regrouping},
combine copies of the same full-germ orbit, and apply
Theorem~\ref{thm:calculus-hom} to the projection onto $K$.  We obtain a finite
orthogonal sum of full-germ bundles, with exactly the intertwiners
in Equation~\eqref{eq:calculus-hom}.  In particular, if its fibre
representations are $\rho_d$, its actual dimensions are
\begin{equation}\label{eq:calculus-summand-dimensions}
 \dim_M K=\sum_d\dim_{\mathbb C}(\rho_d)[\Lambda:B_d],
 \qquad
 \dim K_N=\sum_d\dim_{\mathbb C}(\rho_d)[\Gamma:A_d].
\end{equation}
Both are positive integers for $K\ne0$.  This conclusion uses the
description of its projection; the dimensions of a containing model
alone would not imply it.

\begin{remark}[Why a presenting graph must be enlarged]
Let $M=N=L_\mu(\Gamma)$, with arbitrary $\mu$, and suppose
$A<\Gamma$ has finite index $k>1$.  Present the identity germ $d$
by $\id_A:A\to A$, with the one-dimensional trivial fibre.
The multiplier on every diagonal pair is
$\mu(a,b)\overline{\mu(a,b)}=1$, so this is a model for any source
twist.  Its full stabilizer is the diagonal copy of $\Gamma$.
Regrouping therefore replaces the presenting fibre by the permutation
representation on $\ell^2(\Gamma/A)$.

The constant vectors form an invariant line.  Its orthogonal projection
has every matrix entry equal to $1/k$, and selects
$H(d,1)=L^2(M)$ from the model $E(A,A,\id_A,1)$.
The selected module has dimensions $(1,1)$, whereas the containing model
has dimensions $(k,k)$.  On each of the $k$ original one-dimensional
fibres above $d$, the diagonal block of the selecting projection is $1/k$, which
is not a projection.  Grouping those fibres at the full stabilizer is
what exposes the invariant line and its actual rank.
\end{remark}

\subsection{Conjugation}

On the conjugate correspondence $\overline{\mathcal H(C,\gamma)}$
the endpoint groups are swapped.  Its subgroup and base fibre are
\[
 C^{\mathrm{op}}=\{(h,g):(g,h)\in C\},\qquad
 \bar\gamma(h,g)\bar\xi=\overline{\gamma(g,h)\xi}.
\]
The multiplier is
$\omega\boxtimes\bar\mu$, since
$\bar\gamma$ conjugates $\alpha$ and swaps its two variables.
The label $x$ of a full-germ bundle becomes the inverse germ $x^{-1}$:
indeed
\[
 \bigl(\Ad(h)x\Ad(g^{-1})\bigr)^{-1}
 =\Ad(g)x^{-1}\Ad(h^{-1}).
\]
With swapped coset representatives, the transport is explicitly
\begin{equation}\label{eq:calculus-conjugate-transport}
 U^{\bar X}_{(h,g)}(\bar x)
   =\overline{U^X_{(g,h)}(x)}.
\end{equation}
Thus it includes the conjugate of the scalar ratio in
Equation~\eqref{eq:input-A1}.  At the module level this equality is
the defining identity
$v_h\bar\xi u_g^*=\overline{u_g\xi v_h^*}$.
For an original graph $\delta:A\to B$, the new graph is
$\delta^{-1}:B\to A$ and its fibre representation is
$b\mapsto\overline{\sigma(\delta^{-1}(b))}$.
Fibre intertwiners are conjugated by the same rule; their adjoints and
norms are preserved.  Proposition~\ref{prop:calculus-dimensions}
also shows directly that the left and right dimensions are exchanged.

\subsection{Balanced fusion and its scalar transport}

Let $P=L_\nu(\Pi)$.  Consider a bundle $\mathcal H_X$ from
$(\Gamma,\mu)$ to $(\Lambda,\omega)$ and a bundle $\mathcal H_Y$
from $(\Lambda,\omega)$ to $(\Pi,\nu)$, each constructed as above.
Write $V_x$ and $W_y$ for their fibres and
$\alpha'=\omega\boxtimes\bar\nu$.  Finite direct sums are handled
by distribution, so we may first take one transitive bundle at each
endpoint.  Define
\begin{equation}\label{eq:calculus-balanced-labels}
 Z=(X\times Y)/\Lambda,
 \qquad h(x,y)=\bigl((1,h)x,(h,1)y\bigr).
\end{equation}
Identify the tensor fibres along this action using
\begin{equation}\label{eq:calculus-middle-transport}
 J_h(x,y)=U^X_{(1,h)}(x)\otimes U^Y_{(h,1)}(y).
\end{equation}
These identifications are genuine unitary group transports, because
their multiplier is
$\overline{\omega(h_1,h_2)}\omega(h_1,h_2)=1$.
The action on labels is free, as already its action on $X$ is free.
Each quotient fibre $V_z$ is consequently unambiguously a tensor
fibre after choosing one representative of $z$.  Put
$\mathcal H_Z=\bigoplus_{z\in Z}V_z$, with counting norm at any
such representatives.  Changing representatives gives a unitary
identification by Equation~\eqref{eq:calculus-middle-transport}.

The outer transport before passage to the quotient is
$U^X_{(g,1)}\otimes U^Y_{(1,k)}$.  It commutes exactly with every
$J_h$: each mixed-coordinate factor of either product cocycle is
one.  Its multiplier is $\mu\boxtimes\bar\nu$.
For an explicit formula choose a representative $(x,y)$ of $z$ and
representatives for all quotient labels.  There is a unique $h$ such
that
\[
 \bigl((g,h)x,(h,k)y\bigr)
\]
is the chosen representative of $(g,k)z$.  With
$s=(g,h)$, $t=(h,k)$, $c=c_X(s,x)$ and $b=c_Y(t,y)$, the quotient
transport is
\begin{equation}\label{eq:input-A2}
 \frac{\alpha(s,r_x)\alpha'(t,r_y)}
      {\alpha(r_{sx},c)\alpha'(r_{ty},b)}
       \gamma_X(c)\otimes\gamma_Y(b).
\end{equation}
It is the product of the two transports in
Equation~\eqref{eq:input-A1}.  There is no additional factor from
the compensating element $h$, since its lift is the genuine
transport in Equation~\eqref{eq:calculus-middle-transport}.

To attach a germ to a quotient label, first recall the map from graph
cosets to germs.  For $x=(g,h)C_\delta\in X$, put
\[
 d_x=\Ad(h)[\delta]\Ad(g^{-1}),
\]
and define $d_y$ in the same way for $y\in Y$.  These are the
graph-coset-to-germ maps from the earlier regrouping, applied at their
respective endpoints; they need not be injective.  The middle action changes the two associated germs by
\[
 (d_x,d_y)\longmapsto
       (\Ad(h)d_x,\,d_y\Ad(h^{-1})).
\]
Hence $z=[(x,y)]\mapsto d_y d_x$ is well defined.  Distinct quotient
labels can still have the same germ; the next proposition explains how
their fibres are regrouped.

\begin{proposition}[Fusion as a counting tensor-fibre bundle]
\label{prop:calculus-fusion}
There is a unitary correspondence identification
\begin{equation}\label{eq:calculus-fusion-unitary}
 \mathcal H_X\mathbin{\bar\otimes}_{N}\mathcal H_Y
 \cong\mathcal H_Z
\end{equation}
under which outer transport is Equation~\eqref{eq:input-A2}.
The space $Z$ has finitely many outer-group orbits.  For each
$z=[(x,y)]$, its stabilizer is a finite-index subgroup-isomorphism
graph with germ $d_y d_x$.
Regrouping these orbits over full stabilizers by
Equation~\eqref{eq:calculus-regrouping} gives the full-germ
decomposition of the fusion.
\end{proposition}

\begin{proof}
First we check the Hilbert norm and the relative product explicitly.
Every $\xi\in V_x$ is right bounded: the map
\[
 R_\xi:L^2(N)\longrightarrow\mathcal H_X,
 \qquad \widehat n\longmapsto\xi n
\]
initially defined on the group algebra extends with norm $\|\xi\|$.
Indeed the vectors $\xi v_h$ occupy distinct fibres as $h$ varies,
and all have norm $\|\xi\|$.  Choose one $x_i$ in each right
$\Lambda$-orbit of $X$, and an orthonormal basis $(e_{ij})_j$ of
$V_{x_i}$.  Orthogonality of all their group translates gives
\[
 R_{e_{ij}}^*R_{e_{i'j'}}=\delta_{ii'}\delta_{jj'}1.
\]
Their ranges exhaust $\mathcal H_X$, because the translates exhaust
every labelled fibre.  They therefore form an orthonormal free
right $N$ basis, with finitely many elements.

The defining inner product of the relative tensor product now gives
\begin{equation}\label{eq:calculus-relative-inner-product}
 \left\langle e_{ij}\otimes_N\eta,
                 e_{i'j'}\otimes_N\eta'\right\rangle
 =\delta_{ii'}\delta_{jj'}\langle\eta,\eta'\rangle.
\end{equation}
Thus that product is unitarily
$\bigoplus_i(V_{x_i}\otimes\mathcal H_Y)$.
Every label in $Z$ has a unique representative with first coordinate
one of the fixed $x_i$: freeness of the right action on $X$ gives
uniqueness of the required balancing element.  Consequently the
counting tensor-fibre Hilbert sum $\mathcal H_Z$ is exactly this
same orthogonal sum.  Its map to the relative product sends
$\xi\otimes\eta\in V_x\otimes W_y$ to
$\xi\otimes_N\eta$.  Equation~\eqref{eq:calculus-relative-inner-product}
proves that it is isometric and onto.  It is independent of the
representative, since
\begin{equation}\label{eq:calculus-balanced-vector}
 (\xi v_h^*)\otimes_N(v_h\eta)
 =\xi\otimes_N(v_h^*v_h\eta)
 =\xi\otimes_N\eta.
\end{equation}
Applying the outer generators to these vectors proves
Equation~\eqref{eq:input-A2} at the correspondence level.

For finiteness and the stabilizers we give the explicit double-coset
calculation, which also specifies the composed germ and its fibre
representation.  Write the two graph maps as
\[
 \delta:A\longrightarrow B\leq\Lambda,
 \qquad \varepsilon:C\longrightarrow F\leq\Pi,
\]
with $A\leq\Gamma$ and $C\leq\Lambda$ of finite index, and
take $X=(\Gamma\times\Lambda)/C_\delta$,
$Y=(\Lambda\times\Pi)/C_\varepsilon$.
Choose $m$ in the finite double-coset set $B\backslash\Lambda/C$.
Represent its outer orbit by
\[
 z_m=[(C_\delta,(m,1)C_\varepsilon)].
\]
Outer translation and balancing send an arbitrary pair to the form
\[
 ((g,h)C_\delta,(\ell,k)C_\varepsilon)
 \longmapsto(C_\delta,(h^{-1}\ell,1)C_\varepsilon).
\]
Changing the initial coset representatives replaces $h^{-1}\ell$
by $\delta(a)^{-1}h^{-1}\ell c$ for $a\in A$, $c\in C$.
Hence these outer orbits are indexed exactly by
$B\backslash\Lambda/C$.

An outer element $(a,p)$ stabilizes $z_m$ precisely when some $h$
satisfies
\[
 (a,h)\in C_\delta,
 \qquad (m^{-1}hm,p)\in C_\varepsilon.
\]
Its stabilizer is consequently the graph of
\begin{equation}\label{eq:calculus-composed-graph}
 \begin{split}
 A_m&=\delta^{-1}(B\cap mCm^{-1}),\\
 \theta_m&=\varepsilon\circ\Ad(m^{-1})\circ\delta:A_m\longrightarrow F_m,\\
 F_m&=\varepsilon(C\cap m^{-1}Bm).
 \end{split}
\end{equation}
Both projections have finite index.  In particular
\begin{equation}\label{eq:calculus-double-coset-labels}
 Z\cong\bigsqcup_{m\in B\backslash\Lambda/C}
             (\Gamma\times\Pi)/C_{\theta_m}.
\end{equation}
Here $C_{\theta_m}$ denotes the graph of the displayed actual map,
not necessarily its full-germ stabilizer.
At $z_m$ the composite germ is
$[\varepsilon\Ad(m^{-1})\delta]=[\theta_m]$.
Its graph lies at finite index in its full stabilizer by
Lemma~\ref{lem:calculus-germ-test}.  This proves all the stated
regrouping assertions.
\end{proof}

For completeness, the last proof yields an explicit finite
double-coset formula without unspecified stabilizer representations.
Let $\sigma$ and $\tau$ be the original fibre representations on
$A$ and $C$.  In the fibre at $z_m$ choose $r_{(m,1)C_\varepsilon}=(m,1)$.
For $a\in A_m$, write $b=\delta(a)$ and $c=m^{-1}bm\in C$.
Equation~\eqref{eq:input-A2}, with compensating element $h=b$, gives
\begin{equation}\label{eq:calculus-double-coset-fibre}
 \gamma_m(a)=
 q_m(b)\,\sigma(a)\otimes\tau(m^{-1}bm),
 \qquad
 q_m(b)=\frac{\omega(b,m)}{\omega(m,m^{-1}bm)}.
\end{equation}
The correction has the exact coboundary
\begin{equation}\label{eq:calculus-conjugation-coboundary}
 \frac{q_m(b)q_m(b')}{q_m(bb')}
 =\frac{\omega(b,b')}
 {\omega(m^{-1}bm,m^{-1}b'm)}.
\end{equation}
To verify this, the twisted group-unit identity
$v_bv_m=q_m(b)v_mv_{m^{-1}bm}$ can be multiplied successively for
$b,b'$ and compared with the identity for $bb'$.
Combining Equation~\eqref{eq:calculus-conjugation-coboundary} with
the multipliers of $\sigma$ and $\tau$ shows that $\gamma_m$ has
multiplier $\mu|_{A_m}/\theta_m^*\nu$.  Hence
\begin{equation}\label{eq:calculus-double-coset-fusion}
 E(A,B,\delta,\sigma)\mathbin{\bar\otimes}_N
 E(C,F,\varepsilon,\tau)
 \cong\bigoplus_{m\in B\backslash\Lambda/C}
      E(A_m,F_m,\theta_m,\gamma_m).
\end{equation}
The appropriate endpoint algebras in the second model and in each
summand are understood from their displayed groups.
One then applies finite projective induction to the full stabilizers
and combines any equal full-germ orbits.  Formula
\eqref{eq:calculus-double-coset-fibre} exhibits explicitly why replacing
the twisted fusion rule by its ordinary version would lose a phase.

The indices in Equation~\eqref{eq:calculus-composed-graph} satisfy
\[
 [\Gamma:A_m]=[\Gamma:A][B:B\cap mCm^{-1}],
 \qquad
 [\Pi:F_m]=[\Pi:F][C:C\cap m^{-1}Bm].
\]
Counting the left $B$-orbits in $\Lambda/C$ and the right $C$-orbits
in $B\backslash\Lambda$ gives, respectively,
\[
 \sum_m[B:B\cap mCm^{-1}]=[\Lambda:C],
 \qquad
 \sum_m[C:C\cap m^{-1}Bm]=[\Lambda:B].
\]
Thus both coordinate groups have finitely many orbits in the fused
label space.  These formulas also verify, directly from the models
and their counting norms, the usual multiplicativity of the left
dimensions and of the right dimensions under this fusion.

\subsection{Maps, associativity, and completion of the calculus}

The preceding identifications respect the entire operator calculus,
not only the isomorphism classes of objects.  For the map calculation,
first regroup the inputs into full-germ bundles.
Theorem~\ref{thm:calculus-hom} then represents every bounded intertwiner
by an equivariant family of fibre maps.  An equivariant family
$T_x:V_x\to V'_x$ and an equivariant family $S_y:W_y\to W'_y$
give the fibre maps $T_x\otimes S_y$.  They commute with the
middle identifications in Equation~\eqref{eq:calculus-middle-transport},
so descend to the quotient fibres.  Under the unitary of
Proposition~\ref{prop:calculus-fusion}, they are precisely
$T\mathbin{\bar\otimes}_N S$, as is seen on
$\xi\otimes_N\eta$.  Their norms are bounded by $\|T\|\|S\|$;
adjoints and compositions agree fibrewise.  Regrouping by finite
projective induction is a regrouping of these same Hilbert fibres
and maps.  Conjugation has the map formula already given in
Equation~\eqref{eq:calculus-conjugate-transport}.

For three bundles, both orders of fusion use the quotient of the
triple label space by the two middle groups.  Their actions commute:
on the shared middle bundle they are its commuting coordinate
actions, and their projective factors cancel with those in the
adjacent bundles.  The fibre identification between the two orders
is the ordinary unitary
\[
 (\xi\otimes\eta)\otimes\zeta
       \longmapsto\xi\otimes(\eta\otimes\zeta).
\]
The counting norms and Equation~\eqref{eq:calculus-relative-inner-product}
show that it is also the canonical associativity unitary for the
relative products.  All scalar identifications are those of
Equation~\eqref{eq:calculus-middle-transport}; hence this assertion
requires no change of cocycle.  For any further number of factors
these maps agree on elementary tensors, which are total, so the
associativity identifications are coherent.  The unit correspondence
$L^2(M)$ is the identity-germ model with its one-dimensional trivial
graph representation; its unit maps are ordinary left and right
multiplication on the corresponding relative tensors.

Conjugation reverses a balanced product.  The unitary
\[
 \overline{\mathcal H_X\mathbin{\bar\otimes}_N\mathcal H_Y}
 \longrightarrow
 \overline{\mathcal H_Y}\mathbin{\bar\otimes}_N
       \overline{\mathcal H_X},
 \qquad
 \overline{\xi\otimes_N\eta}
       \longmapsto\bar\eta\otimes_N\bar\xi
\]
is well defined on tensors of vectors in labelled fibres: the
simultaneous middle transport
becomes the same transport on the swapped conjugate fibres, with
both scalar factors conjugated and their product still one.
The counting tensor-fibre model of
Proposition~\ref{prop:calculus-fusion} proves preservation of the norm:
the map bijects the quotient labels and is unitary on each conjugate
tensor fibre.  These tensors are total.  It intertwines the swapped
endpoint actions and the conjugated fibre maps.  Thus the conjugation,
unit, and associativity identifications agree with those of the
correspondence category at the level of actual operators.

We now apply the geometric input. Proposition~\ref{prop:recovery-exhaustion}
places every bifinite correspondence with source in $\mathscr K$ inside
a finite sum of graph models. The preceding full-germ calculation resolves
its selecting projection into fibre projections. Together with the operation
formulas just proved, this establishes every assertion of
Theorem~\ref{thm:correspondences}, including
all bounded intertwiners, actual summands and their dimensions, and
the specified conjugation and fusion maps.  Throughout the section
only the matched-subgroup ratio is required to have a
finite-dimensional projective realization.  None of the calculations
imposes a finite-image condition on that representation or a
finite-type condition on either individual endpoint cocycle.
